\label{chap:csz}

\begin{definition}[CSZ system]
  \label{def:csz}
  \lean{CantorSzemeredi.MPS.Factor.CSZ, CantorSzemeredi.MPS.multiCorr}
  \leanok
  \uses{def:mps, def:cantor, def:uniform-liminf}
  A measure preserving system $(X,\Abb,m,T)$ is \emph{Cantor--Szemer\'edi}, or \emph{CSZ}, if
  for all measurable $A \subset X$ with $m(A) > 0$, all integers $t \geq 1$, all integers
  $i \geq 0$, and all Cantor sets $\Cantor = \Cantor_{b,D}$ with $0 \in D$,
  \begin{equation}
    \label{eq:csz}
    \unifliminf \ \unifavg \int_X \prod_{s=0}^{t} T^{s b^i k_n} \one_A \, dm > 0 .
  \end{equation}
  A factor $\Dbb \subset \Abb$ is CSZ if $(X,\Dbb,m,T)$ is.
\end{definition}

\begin{remark}[Lean form]
  \leanok
  In Lean the integral in \eqref{eq:csz} is written as the measure of an intersection,
  $\texttt{multiCorr A t d} = m\bigl(\bigcap_{s=0}^{t} T^{-sd}A\bigr)$ with $d = b^ik_n$; the two
  agree because a product of indicators is the indicator of the intersection, and the integral
  of an indicator is the measure.  \texttt{Factor.CSZ D} quantifies over $\Cantor$ (as a
  \texttt{CantorData}), over $A$ with \texttt{MeasurableSet[D.\ensuremath{\sigma}] A} and
  $0 < m(A)$, over $t \geq 1$ and over $i$, and asserts
  $0 < \texttt{uniformLiminf (fun n => multiCorr A t (base\^{}i * enum n))}$.
\end{remark}

\begin{lemma}[Sets versus non-negative functions]
  \label{lem:csz-functions}
  \lean{CantorSzemeredi.MPS.csz_functions_iff, CantorSzemeredi.MPS.csz_functions_of_csz}
  \leanok
  \uses{def:csz}
  Let $\Dbb$ be a factor, $\Cantor$ a Cantor set, $t \geq 1$ and $i \geq 0$.  The following
  are equivalent.
  \begin{enumerate}
    \item[(a)] \eqref{eq:csz} holds for every $A \in \Dbb$ with $m(A) > 0$.
    \item[(b)] For every bounded $\Dbb$-measurable $f \geq 0$ which is not a.e.\ zero,
      \[
        \unifliminf \unifavg \int_X \prod_{s=0}^{t} T^{sb^ik_n} f\,dm > 0 .
      \]
  \end{enumerate}
\end{lemma}

\begin{proof}
  \uses{def:csz}
  (b) $\Rightarrow$ (a): take $f = \one_A$, which is bounded, $\Dbb$-measurable, non-negative,
  and not a.e.\ zero since $m(A) > 0$.

  (a) $\Rightarrow$ (b): let $f$ be as in (b).

  \emph{A level set of positive measure.}  There is $c > 0$ with $m(\{f > c\}) > 0$: otherwise
  $m(\{f > 1/k\}) = 0$ for every $k \geq 1$, so $\{f > 0\} = \bigcup_k\{f > 1/k\}$ is null and
  $f = 0$ a.e.  Fix such $c$ and put $A \coloneqq \{f > c\} \in \Dbb$.

  \emph{Pointwise domination.}  $f \geq c\,\one_A$ pointwise: on $A$, $f > c = c\one_A$; off
  $A$, $f \geq 0 = c\one_A$.  Composing with $T^{sb^ik_n}$ preserves this, and all factors are
  non-negative, so for every $n$,
  \[
    \prod_{s=0}^{t} T^{sb^ik_n}f \geq \prod_{s=0}^{t} T^{sb^ik_n}(c\one_A) = c^{t+1}\prod_{s=0}^{t} T^{sb^ik_n}\one_A
    \quad\text{pointwise, hence}\quad
    \int_X \prod_{s} T^{sb^ik_n}f\,dm \geq c^{t+1} \int_X \prod_s T^{sb^ik_n}\one_A\,dm .
  \]
  (Monotonicity of a product of non-negative factors in each factor: induct on the number of
  factors, or use \texttt{Finset.prod\_le\_prod}.)

  \emph{The uniform liminf.}  Averaging over a window preserves the inequality, and
  \texttt{uniformLiminf} is monotone and commutes with multiplication by the positive constant
  $c^{t+1}$.  So the uniform liminf in (b) is at least $c^{t+1}$ times the one in \eqref{eq:csz}
  for $A$, which is positive by (a).
\end{proof}

\begin{remark}[On the dilation parameter $i$]
  \label{rem:dilation}
  \uses{cor:self-similar}
  Two points, both needed in Lean.
  \begin{enumerate}
    \item Since $b^i k_n = k_{\delta^i n}$ (Corollary~\ref{cor:self-similar}), the condition for
      a given $i$ is the case $i = 0$ with the index restricted to the progression
      $\delta^i \NNz$; equivalently, it is the case $i=0$ for the system
      $(X,\Abb,m,T^{b^i})$ and the \emph{same} Cantor set $\Cantor$.
    \item That equivalence is \emph{not} the statement that a CSZ system $(X,\Abb,m,T)$ gives
      rise to a CSZ system $(X,\Abb,m,T^{b^i})$.  The CSZ property quantifies over all Cantor
      sets, and for a second Cantor set the available dilations are powers of \emph{its} base,
      not of $b$.  In Lean the single instance, for one Cantor set at a time, is what must be
      used; stating the stronger (false) version will break the compact-extension proof.
  \end{enumerate}
  Carrying $i$ is not optional: the compact extension argument uses the hypothesis on the factor
  at a dilation determined inside that proof, and returns the conclusion for the extension at the
  dilation it was handed.  It is the self-similarity $b^i \Cantor \subset \Cantor$ that makes
  this work.  Remark~\ref{rem:compact-organization} spells out the resulting organization.
\end{remark}

\begin{lemma}[The trivial factor is CSZ]
  \label{lem:trivial-csz}
  \lean{CantorSzemeredi.MPS.Factor.trivial_CSZ}
  \leanok
  \uses{def:csz}
  The factor consisting of the null and conull sets is CSZ.
\end{lemma}

\begin{proof}
  \leanok
  \uses{def:csz}
  If $A$ lies in it and $m(A) > 0$ then $A$ is conull, so every integral in \eqref{eq:csz}
  equals $1$.
\end{proof}

\begin{lemma}[Extension principles]
  \label{lem:extensions}
  \lean{CantorSzemeredi.MPS.csz_sets_of_compactExtension}
  \uses{def:csz, def:cwm, def:compact-extension, lem:wm-extension, thm:compact-extension}
  Let $\Dbb \subset \Ebb$ be factors of $(X,\Abb,m,T)$ with $\Dbb$ CSZ.
  \begin{enumerate}
    \item[(W)] If $\Ebb$ is a weakly mixing extension of $\Dbb$, then $\Ebb$ is CSZ.
    \item[(K)] If $\Ebb$ is a compact extension of $\Dbb$, then $\Ebb$ is CSZ.
  \end{enumerate}
  % Formal status: (K) is \texttt{csz\_sets\_of\_compactExtension}, proved in the set form;
  % (W) is not yet formalized, since it rests on \texttt{lem:wm-extension} and hence on
  % \texttt{lem:wm-main}.  This lemma is therefore not yet marked as formalized.
\end{lemma}

\begin{proof}
  \uses{lem:wm-extension, thm:compact-extension}
  (W) is Lemma~\ref{lem:wm-extension}; (K) is Theorem~\ref{thm:compact-extension}, in its set
  form $f = \one_A$, $A \in \Ebb$.
\end{proof}

\begin{lemma}[Compact systems are compact extensions of the trivial factor]
  \label{lem:compact-system-extension}
  \lean{CantorSzemeredi.MPS.compactExtension_of_compactSystem}
  \leanok
  \uses{def:compact-system, def:compact-extension}
  If $(X,\Abb,m,T)$ is a compact system then $\Abb$ is a compact extension of the trivial
  factor.
\end{lemma}

\begin{proof}
  \uses{def:compact-system, def:compact-extension}
  For the trivial factor $\Dbb$, $\cnorm{u}(x) = \norm{u}_{L^2(m)}$ for a.e.\ $x$ (the
  conditional expectation onto $\Dbb$ is the integral).  Let $f \in L^2(X)$ and $\epsilon > 0$.
  By Definition~\ref{def:compact-system} there are $r$ and $g_1,\dots,g_r$ with: for every
  $n \in \ZZ$ there is $q$ with $\norm{f\circ T^n - g_q}_2 \leq \epsilon/2$.  The $g_q$ need not
  be measurable or square integrable, so replace them by orbit elements: for each $q$ for
  which some $n$ has $\norm{f\circ T^n - g_q}_2 \leq \epsilon/2$, pick one such $n_q$ and set
  $h_q \coloneqq f \circ T^{n_q}$; discard the other $q$.  Each $h_q$ is measurable and in $L^2$
  (as $T^{n_q}$ preserves $m$), so $h_q \in L^2(\Abb)$ (\texttt{MemL2} for the top factor).
  For every $n \in \NN$ choose $q$ with $\norm{f\circ T^n - g_q}_2 \leq \epsilon/2$; then
  $\norm{f\circ T^n - h_q}_2 \leq \norm{f\circ T^n - g_q}_2 + \norm{g_q - f\circ T^{n_q}}_2 \leq \epsilon$,
  and this is $\cnorm{f \circ T^n - h_q}(x) \leq \epsilon$ for a.e.\ $x$ (a constant).  Using
  $\epsilon/2$ in place of $\epsilon$ gives the strict inequality required by
  \texttt{CondCompact}.  So every $f \in L^2(X)$ is conditionally compact over the trivial
  factor, and in particular the conditionally compact functions are dense in $L^2(\Abb)$.
\end{proof}

\begin{corollary}
  \label{cor:wm-compact-are-csz}
  \lean{CantorSzemeredi.MPS.CSZ_of_compactSystem, CantorSzemeredi.MPS.CSZ_of_weaklyMixing}
  \leanok
  \uses{lem:extensions, lem:trivial-csz, lem:wm-sets-to-functions, lem:compact-system-extension}
  Weakly mixing systems and compact systems are CSZ.
\end{corollary}

\begin{proof}
  \uses{lem:extensions, lem:trivial-csz, lem:wm-sets-to-functions, lem:compact-system-extension}
  A weakly mixing system is a weakly mixing extension of the trivial factor by
  Lemma~\ref{lem:wm-sets-to-functions}, and a compact system is a compact extension of it by
  Lemma~\ref{lem:compact-system-extension}; apply Lemma~\ref{lem:trivial-csz} and
  Lemma~\ref{lem:extensions}.  Note that CSZ for the top factor $\Abb$ is CSZ for the system.
\end{proof}
